% =====================================================================
% Conecthus - Entrega do Teste Técnico (Instituto Conecthus)
% Author: Arao Melo Lopes - 21/09/2026
% Compilar: pdflatex conecthus_entrega.tex
% =====================================================================
\documentclass[11pt,a4paper]{article}

\usepackage[T1]{fontenc}
\usepackage[brazilian]{babel}
\usepackage{lmodern}
\usepackage[margin=2.2cm]{geometry}
\usepackage{booktabs}
\usepackage{tabularx}
\usepackage{array}
\usepackage{enumitem}
\usepackage[colorlinks=true,linkcolor=black,urlcolor=blue]{hyperref}
\usepackage{xcolor}
\usepackage{parskip}

\setlength{\parindent}{0pt}
\urlstyle{same}

\begin{document}

% ---------------------------------------------------------------------
% CAPA
% ---------------------------------------------------------------------
\begin{titlepage}
  \centering
  {\LARGE\bfseries Conecthus --- Task Manager\par}
  \vspace{1.2cm}
  {\large Teste Técnico\par}
  \vspace{0.3cm}
  {\large Instituto Conecthus\par}
  \vspace{1.5cm}
  {\large Aarão Melo Lopes\par}
  \vspace{4cm}

  \begin{tabular}{@{}ll@{}}
    \toprule
    Data da entrega & 21 de setembro de 2026 \\[3pt]
    Repositório & \url{https://github.com/aaraomelo/conecthus} \\[3pt]
    Aplicação Web & \url{https://conecthus.patriatechnology.com} \\[3pt]
    Swagger (API) & \url{https://conecthus.patriatechnology.com/api/docs} \\
    \bottomrule
  \end{tabular}

  \vfill
  {\footnotesize Documento de entrega --- aplicação, testes e publicação.}
\end{titlepage}

% ---------------------------------------------------------------------
% 1. RESUMO DA SOLUCAO
% ---------------------------------------------------------------------
\section{Resumo da solução}

O \textbf{Conecthus} é um gerenciador de tarefas full-stack (CRUD) com autenticação
JWT, validação no cliente e no servidor, cache, busca textual e notificações em
tempo real via MQTT. A aplicação é entregue como uma única origem (\emph{single-origin}):
a SPA Web, servida por nginx, faz proxy de \texttt{/api} (REST) e \texttt{/mqtt}
(WebSocket) para o backend e o broker MQTT.

\begin{itemize}[leftmargin=1.4em,itemsep=2pt]
  \item \textbf{Web:} React 19 + Vite + TypeScript + Redux Toolkit, responsivo, com
        React Hook Form + Zod para validação de formulários.
  \item \textbf{Backend:} NestJS (API REST) + Prisma sobre PostgreSQL, Redis para cache,
        MQTT para notificações e Swagger para documentação da API.
  \item \textbf{Mobile:} React Native (Expo SDK 57), consumindo a mesma API e os
        mesmos tópicos MQTT.
  \item \textbf{Infraestrutura:} Docker Compose (PostgreSQL 16, Redis 7, Mosquitto 2.0.18,
        backend, frontend), com healthchecks e deploy contínuo via GitHub Actions.
\end{itemize}

% ---------------------------------------------------------------------
% 2. STACK
% ---------------------------------------------------------------------
\section{Stack}

\begin{tabularx}{\linewidth}{@{}l >{\raggedright\arraybackslash}X@{}}
  \toprule
  \textbf{Componente} & \textbf{Tecnologias} \\
  \midrule
  Web
    & React 19, Vite 8, TypeScript, Redux Toolkit, React Router 7,
      React Hook Form + Zod (validação), MQTT.js; testes com Vitest + Testing Library,
      lint com oxlint \\
  Backend
    & NestJS 12, Prisma 7, class-validator (ValidationPipe), JWT
      (@nestjs/jwt), Swagger (@nestjs/swagger), cache-manager + ioredis, mqtt.js \\
  Banco de dados & PostgreSQL 16 com Prisma 7 (schema + migrations) \\
  Cache & Redis 7 (TTL 60s; invalidação por usuário após mutações) \\
  Mensageria & Eclipse Mosquitto 2.0.18 --- notificações em tempo real
      em \texttt{notifications/\{userId\}} (MQTT 1883; WebSocket 8083) \\
  Mobile & React Native 0.86 + Expo SDK 57, React Navigation, axios,
      AsyncStorage, MQTT via WebSocket; testes com jest-expo \\
  Infraestrutura & Docker Compose (Postgres, Redis, Mosquitto, backend, frontend/nginx);
      deploy contínuo no GitHub Actions \\
  \bottomrule
\end{tabularx}

% ---------------------------------------------------------------------
% 3. FUNCIONALIDADES IMPLEMENTADAS
% ---------------------------------------------------------------------
\section{Funcionalidades implementadas}

\begin{itemize}[leftmargin=1.4em,itemsep=2pt]
  \item \textbf{Autenticação:} cadastro e login com JWT (senha com bcrypt);
        guard global no backend com opt-out apenas em \texttt{auth/*}.
  \item \textbf{Proteção de rotas:} rotas privadas/públicas na Web
        (\texttt{ProtectedRoute}/\texttt{GuestRoute}) e navegação condicional no Mobile.
  \item \textbf{CRUD de tarefas:} criar, listar, ver, editar, concluir e excluir,
        via REST + Prisma, com isolamento por usuário (\texttt{userId} no \texttt{WHERE}).
  \item \textbf{Filtros:} por \texttt{status} e período (\texttt{dueDate}).
  \item \textbf{Paginação:} cursor/offset com parâmetros \texttt{page}/\texttt{pageSize}.
  \item \textbf{Busca:} textual em \texttt{title}/\texttt{description} via
        \texttt{ILIKE} (decisão documentada no repositório).
  \item \textbf{Cache Redis:} listas e perfis com TTL 60s e invalidação
        por usuário após qualquer mutação.
  \item \textbf{MQTT:} notificações em tempo real em \texttt{notifications/\{userId\}},
        consumidas na Web (toasts) e no Mobile (overlay).
  \item \textbf{Web:} SPA responsiva com formulários validados (React Hook Form + Zod).
  \item \textbf{Mobile:} Expo SDK 57 --- telas de autenticação, lista de tarefas com
        filtros/busca/paginação e notificações.
  \item \textbf{Swagger:} documentação da API em \texttt{GET /api/docs}.
  \item \textbf{Docker:} ambiente completo com \texttt{docker compose up --build -d}
        e healthchecks para todos os serviços.
\end{itemize}

% ---------------------------------------------------------------------
% 4. QUALIDADE E TESTES
% ---------------------------------------------------------------------
\section{Qualidade e testes}

\begin{tabularx}{\linewidth}{@{}l l >{\raggedright\arraybackslash}X@{}}
  \toprule
  \textbf{Item} & \textbf{Resultado} & \textbf{Observação} \\
  \midrule
  Testes unitários Web (Vitest) & 77/77 aprovados
    & 12 arquivos de teste; inclui \texttt{TaskFormPage} 3/3 \\
  Testes unitários Backend (Vitest) & 26/26 aprovados
    & 6 arquivos de teste \\
  Testes com \texttt{skip}/\texttt{only}/\texttt{todo} & Nenhum
    & Verificado por varredura no código-fonte \\
  Lint oxlint --- Web & 0 erros; 0 warnings
    & \texttt{npm run lint} \\
  Lint oxlint --- Backend & 0 erros; 0 warnings
    & \texttt{npm run lint} \\
  Build Web & OK (passa no CI)
    & \texttt{tsc -b \&\& vite build} \\
  Build Backend & OK (passa no CI)
    & \texttt{nest build} \\
  Cobertura & Thresholds 70\% configurados (Web e Mobile)
    & Configuração de threshold; não é medição desta entrega \\
  Integração (e2e) & 11 casos no repositório
    & 11 casos (auth, CRUD, cache, MQTT e ownership);
      execução via \texttt{npm run test:e2e} (requer infra Docker) \\
  \bottomrule
\end{tabularx}

\medskip
\emph{Nota transparente:} o pipeline do GitHub Actions \textbf{valida builds} (backend e
frontend) e \texttt{docker compose config}, mas \textbf{não executa as suítes de teste};
esses números refletem a execução local documentada acima.

% ---------------------------------------------------------------------
% 5. PUBLICACAO
% ---------------------------------------------------------------------
\section{Publicação}

\begin{tabularx}{\linewidth}{@{}l >{\raggedright\arraybackslash}X@{}}
  \toprule
  \textbf{Item} & \textbf{Estado} \\
  \midrule
  Aplicação Web & \url{https://conecthus.patriatechnology.com} --- HTTP 200
    (cadastro e login auto-atendidos; crie uma conta na primeira visita) \\
  Swagger público & \url{https://conecthus.patriatechnology.com/api/docs} --- HTTP 200 \\
  API pública & \texttt{/api/*} responde no domínio público
    (sem token, \texttt{/api/tasks} retorna 401) \\
  MQTT público (WebSocket) & \texttt{wss://conecthus.patriatechnology.com/mqtt}
    (notificações em tempo real) \\
  Docker & Serviços publicados via Docker Compose no servidor de produção \\
  Deploy & GitHub Actions (push para \texttt{main}) com deploy SSH +
    \texttt{docker compose up}; deploy mais recente \textbf{concluído com sucesso}
    (validação de \texttt{/api/docs} e smoke test HTTP 200 no frontend) \\
  \bottomrule
\end{tabularx}

\medskip
\textbf{Para o avaliador:} acessar a aplicação, criar uma conta e usar o fluxo
completo (criar, editar, concluir e excluir tarefas, aplicar filtros e busca). Não há
credenciais de demonstração públicas; o cadastro é livre.

% ---------------------------------------------------------------------
% 6. HISTORICO TECNICO
% ---------------------------------------------------------------------
\section{Histórico técnico (commits finais)}

\begin{tabularx}{\linewidth}{@{}l >{\raggedright\arraybackslash}X@{}}
  \toprule
  \textbf{Commit} & \textbf{Mensagem} \\
  \midrule
  \texttt{c534915} & \texttt{docs: document text search decision} \\
  \texttt{791929c} & \texttt{fix(mqtt): pin mosquitto to websocket-compatible version} \\
  \texttt{7e99df8} & \texttt{feat(web): migrate auth and task forms to react-hook-form + zod} \\
  \bottomrule
\end{tabularx}

\smallskip
\emph{Nota:} na sequência foi publicado o commit \texttt{a4c7243}
(correção de build detectada pela CI --- tipagem de formulário e código morto),
com novo deploy validado com sucesso.

% ---------------------------------------------------------------------
% 7. SUGESTAO DE SCREENSHOTS
% ---------------------------------------------------------------------
\section{Screenshots sugeridos para a apresentação}

\begin{enumerate}[leftmargin=1.4em,itemsep=2pt]
  \item Tela de login e de cadastro com validação de campos (mensagens de erro).
  \item Lista de tarefas com filtros (status/período), busca e paginação.
  \item Formulário de criação/edição de tarefa com validação.
  \item Notificação em tempo real ao criar ou concluir uma tarefa.
  \item Swagger público (\texttt{/api/docs}) exibindo os endpoints documentados.
  \item Terminal com o resumo dos testes: Web 77/77 e Backend 26/26.
  \item Status do deploy mais recente no GitHub Actions (run concluída com sucesso).
\end{enumerate}

\end{document}